
We set out to test whether machine learning research suffers from epistemic lock-in---a self-reinforcing cycle where benchmark concentration compounds over time. Our findings reveal a more nuanced picture: the field co-evolves with its evaluation standards without becoming trapped by them.

Our analysis of 12,600 papers across 21 venue-years yields three principal findings:

\begin{enumerate}
\item \textbf{Benchmark concentration is measurable and systematic.} HHI provides a tractable metric, revealing moderate concentration (range 0.007--0.046) with high correlation to top-5 dataset dominance ($\rho$ = 0.90).
\end{enumerate}

\begin{enumerate}
\item \textbf{Citation networks serve as conduits for benchmark propagation.} Papers that cite each other exhibit substantially higher benchmark overlap (Jaccard 0.318 vs. 0.014, Cohen's d = 1.93).
\end{enumerate}

\begin{enumerate}
\item \textbf{No evidence for temporal lock-in.} The lagged regression coefficient is positive rather than negative ($\beta$ = +1.60, p = 0.098). Granger causality tests are inconclusive due to insufficient data.
\end{enumerate}

The fear of epistemic lock-in assumes that benchmark adoption is a one-way ratchet. Our evidence suggests otherwise: the machine learning community collectively selects evaluation standards, transmits them efficiently through citation networks, but retains the capacity to revise them---co-evolving with its benchmarks rather than being captured by them.

